\uuid{ORaO}
\exo7id{4531}
\titre{exo7 4531}
\auteur{quercia}
\organisation{exo7}
\datecreate{2010-03-14}
\isIndication{false}
\isCorrection{false}
\chapitre{Suite et série de fonctions}
\sousChapitre{Autre}
\cpge{mpsi-mp2i.fonctions-valeurs-reelles-complexes.c, pcsi.techniques-calcul-fonctions.c, ptsi.techniques-calcul-fonctions.c, mpsi-mp2i.limites-continuite.b, pcsi.limites-continuite.b, ptsi.limites-continuite.b, bcpst1.fonctions-reelles-usuelles.c, bcpst1.suites-reelles, mp-mpi.series-entieres.b, pc.series-entieres.d, psi.series-entieres.d, pt.series-entieres.d, mp-mpi.series-entieres.d, pc.series-entieres.c, psi.series-entieres.c, pt.series-entieres.c}
\module{Analyse}
\niveau{L2}
\difficulte{}

\contenu{
\question{
Soit $\zeta(x) = \sum_{n=1}^\infty \frac1{n^x}$ et
$\gamma = \lim_{n\to\infty}\left(\frac 11 + \dots + \frac 1n -\ln(n) \right)$.

Montrer que $\gamma = 1 + \sum_{n=2}^\infty\left(\frac 1n + \ln\left(1-\frac 1n\right)\right)$
puis que $\gamma = 1 - \sum_{k=2}^\infty\frac{\zeta(k)-1}k$.
}
}
